\chapter{Primary Instabilities}

\begin{itemize}
%%%%%%%%%%%%%%%%%%%%
\item Learning Goals\\
%%%%%%%%%%%%%%%%%%%%  
  By the end of this chapter, students should be able to:

  \begin{enumerate}

  \item Describe the physical setup and relevance of RT and KH
    instabilities in astrophysical contexts (e.g., supernova remnants,
    accretion flows).

  \item Linearize the equations of motion for a fluid interface and apply perturbation analysis.

  \item Derive and interpret the dispersion relation for both RT and KH instabilities.

  \item Identify criteria for instability and relate to astrophysical phenomena.

  \item Compare and contrast the behavior and growth rates of RT and KH
    instabilities.

  \item Describe the physical origin of Jeans and Toomre instabilities in self-gravitating media.

  \item Derive and interpret the relevant dispersion relations for
    Jeans and Toomre instabilities.

  \item Explain the roles of pressure, gravity, rotation, and shear in stabilizing/destabilizing a system.

  \item Apply instability criteria (Jeans length, Toomre Q) to various astrophysical contexts.

  \item Connect intuitive ideas (e.g., pressure vs gravity, centrifugal balance) to the mathematics of stability.
   
  \end{enumerate}    
%%%%%%%%%%%%%%%%%
\item Lesson Plan
%%%%%%%%%%%%%%%%%  
  \begin{enumerate}
  \item Introduction \& Motivation (10 min)
    \begin{itemize}
    \item Visuals to show where RT and KH instabilities appear in astrophysics (e.g., SN remnants, solar).
    \item Guiding question: What happens at the interface between fluids with different densities or velocities?
    \end{itemize}
  \item Physical Setup: RT Instability (10 min)
    \begin{itemize}
    \item  Two-fluid interface with gravity. Discuss initial conditions and physical intuition behind RT instability.
    \end{itemize}
  \item Linear Stability Analysis: RT (15 min)
    \begin{itemize}
    \item  Walk through linearizing Euler’s equations
    \item Bernoulli theorem + boundary conditions.
    \item Derive dispersion relation:
    \item Discuss physical meaning: when is $\omega$ imaginary
      (instability)? Dependence on wave number $k$.
    \end{itemize}
  \item {Quick question}
    \begin{itemize}
      \item What happens if $\rho_1=\rho_2$?
      \end{itemize}
    \item Physical Setup: KH Instability (10 min)
    \begin{itemize}
    \item Introduce shear flow across a contact discontinuity.
    \item Emphasize velocity discontinuity and pressure balance
      \end{itemize}

    \item Linear Stability Analysis: KH (15 min)
    \begin{itemize}
      \item Derive the KH dispersion relation:
      \end{itemize}

    \item RT vs KH Summary Discussion (5 min)
    \begin{itemize}
      \item Venn diagram or table comparing: driving force, required
        conditions, typical growth rates.
      \item Ask students to contribute examples from astrophysics.
      \end{itemize}

    \item Wrap-Up \& Exit Question (5 min)
    \begin{itemize}
      \item Assign exit question: “In a stratified atmosphere with
        both shear and gravity, which instability dominates and why?”
      \item Non-linear evolution and simulations.
      \end{itemize}

    \end{enumerate}
\end{itemize}

\section{Rayleigh-Taylor instability}

Consider two incompressible, inviscid fluids.

Fluid 1, with density $\rho_1$, lies above fluid 2, with density
$\rho_2$. Gravity acts downward, so $\v{g}=-g\hatz$. The fluids have a
contact discontinuity at $z=0$, and we consider small perturbations to
the interface, $\eta(x,t)$. 

The fluid is incompressible and at rest, we consider it irrotational
and thus a potential flow, writing $\v{u}_i =\grad{\phi_i}$ for each
fluid. 

Since the flow is potential, it obeys Laplace's equation

\beq
\Laplace{\phi_i}=0
\eeq

\noindent with the boundary conditions $\phi\vert{z=\pm\infty} =  0$. 

At the interface the contact discontinuity condition of same velocity
for both fluids implies that the velocity of the boundary, $u_z =
\frac{\partial \eta}{\partial t}$
must be the velocity of the fluids 

\beq
\frac{\partial \eta}{\partial t} = \frac{\partial \phi_i}{\partial z}. 
\label{eq:contactboundary}
\eeq

Finally we apply Euler's equation

\beq
\frac{\partial \v{u}}{\partial t} + \left(\v{u}\cdot\grad\right)\v{u} = -\frac{1}{\rho}\grad{p} + \v{g} 
\eeq

substituting $\v{u}=\grad{\phi}$, then $\left(\v{u}\cdot\grad\right)\v{u} = \frac{1}{2}\grad{u^2}$
for irrotational flow and setting $\frac{1}{\rho}\grad{p}=\grad{h}$ for
enthalpy, we write almost all terms under a single gradient


\beq
\grad{\left(\frac{\partial \phi}{\partial t} +  \frac{1}{2}u^2 + h\right)} = \v{g} 
\eeq

and integrating in z

\beq
\frac{\partial \phi}{\partial t} +  \frac{1}{2}u^2 + h + gz = {\rm const}
\eeq

which is Bernoulli's theorem. Applied to the boundary, for an unperturbed state, $u=0$ and $z=0$, we
see that the constant has to be the unpertubed enthalpy, $h_0$. 

\beq
\frac{\partial \phi}{\partial t} +  \frac{1}{2}u^2 + h + g\eta = h_0
\eeq

we thus write this in terms of the perturbed enthalphy

\beq
\rho\left(\frac{\partial \phi}{\partial t} +  \frac{1}{2}u^2 + g\eta\right) = -p^\prime
\eeq

and since the pressure is the same for both fluids

\beq
\rho_1\left(\frac{\partial \phi_1}{\partial t} +  \frac{1}{2}u^2 + g\eta\right) = \rho_2\left(\frac{\partial \phi_2}{\partial t} +  \frac{1}{2}u^2 + g\eta\right)
\label{eq:bernoulli-boundary}
\eeq



We now seek solutions for the three functions $\eta$, $\phi_1$ and
$\phi_2$. We apply now Fourier perturbations and consider the ansatz

\beq
\phi_1 = \hat{\phi}_1 Z_1(z) e^{i\left(kx-\omega t\right)}
\eeq

where $\hat{\phi}_1$ is an amplitude.  The ansatz
respects the symmetry of the problem. Because $\phi_1$ fulfills
the Laplace equation, we find

\beq
\frac{\partial^2 Z_1}{\partial z^2} = k^2 Z_1
\eeq

which has solution

\beq
Z_1(z) = C_1e^{kz} + C_2 e^{-kz}  
\eeq

Since for $z=\infty$ it must vanish, we set $C_1=0$. In a similar way, we can make an
ansatz for $\phi_2$ and conclude that its z-dependence can only go as
exp$(kz)$ due the boundary condition at $z=-\infty$.
This therefore leads to the following three equations for a single Fourier mode


\beqn
\eta(x,t) &=& \hat{\eta} e^{i\left(kx-\omega t\right)}\\
\phi_1(x,z,t) &=& \hat{\phi}_1 e^{-kz}e^{i\left(kx-\omega t\right)}\\
\phi_2(x,z,t) &=& \hat{\phi}_2e^{kz}e^{i\left(kx-\omega t\right)}
\eeqn

Applying the perturbation to the contact discontinuity velocity
condition \eq{eq:contactboundary} for each fluid 

\beqn
-i\omega \hat{\eta} &=& -k \hat{\phi}_1  \label{eq:pert1}\\
-i\omega \hat{\eta} &=& k \hat{\phi}_2 \label{eq:pert2}
\eeqn

which yields $\phi_1 = -\phi_2$ and we can drop one of the equations. 

For the last equation, we apply \eq{eq:bernoulli-boundary} to
infinitesimal perturbations in velocity, discarding the $u^2$ term

\beq
\rho_1\left(\frac{\partial \phi_1}{\partial t} +  g\eta\right) = \rho_2\left(\frac{\partial \phi_2}{\partial t} +  g\eta\right)
\eeq

we thus find, applying the pertubation  

\beq
\rho_1\left(-i\omega\hat{\phi}_1 +  g\hat{\eta}\right) = \rho_2\left(-i\omega\hat{\phi}_2 +  g\hat{\eta}\right)
\label{eq:pert3}
\eeq

and substituting $\phi_2$ for $-\phi_1$, 

\beq
\left(\rho_1 -\rho_2\right)g\hat{\eta}  = i\omega\hat{\phi}_1 \left(\rho_1+\rho_2\right)
\eeq

now substituting \eq{eq:pert1}

\beq
\omega^2  = \frac{\left(\rho_2-\rho_1\right)}{\left(\rho_1+\rho_2\right)} gk  
\eeq

We see that for $\rho_1 > \rho_2$, i.e. the denser fluid lies on top, unstable solutions with
$\omega^2 < 0$ exist. This is the so-called Rayleigh-Taylor instability. It is in essence
buoyancy driven and leads to the rise of lighter material in a stratified atmosphere.
The free energy that is tapped here is the potential energy in the gravitational field.
Also notice that for an ideal gas, arbitrary small wavelengths are unstable, and those
modes will also grow fastest.

If on the other hand we have $\rho_1 < \rho_2$, then the interface is stable and will only
oscillate when perturbed.

The quantity 

\beq
{\rm At} = \frac{\left(\rho_2-\rho_1\right)}{\left(\rho_1+\rho_2\right)}
\eeq

is called Atwood number. 

\subsection{Physical interpretation}

Let us examine what happens when a fluid parcel is translated to the new vertical position $z=
\eta(x)$ \cite{Piriz+06}. The pressure in an incompressible
fluid increases linearly with depth, the elements at a deeper
position $\eta > 0$ will feel a pressure greater than $p_0$. The
elements at $\eta < 0$ will feel a pressure less than $p_0$. The pres-
sure increases proportionally to the fluid density so for
$\eta > 0$, it increases more on the side of the fluid with higher
density. The pressure of the fluids on each side of the trans-
lated interface is

\beqn
p_1^\prime &=& p_0+\rho_1 g \eta\\
p_2^\prime &=& p_0+\rho_2 g \eta
\eeqn

In the new position, a pressure difference $\Delta p= \left(\rho_2- \rho_1\right) g\eta$ is
created across the interface, which tends to deform it further.
This pressure difference drives the motion of the interface.
We can describe this motion by means of Newton’s second
law of motion,

\beqn
m\ddot{\eta} &=& \Delta p A\\
&=& \left(\rho_2- \rho_1\right) g\eta A
\label{eq:RT-Newton}
\eeqn

where $A$ is the area of the interface and $m$ is the mass of the
fluids involved in the motion. To calculate this mass we use what
we've found, that the perturbations
decay from the interface as exp$(-kz)$. Therefore the total effective mass
that participates in the motion is the mass contained within
this distance,

\beq
m=m_1+m_2  = \frac{A}{k}\left(\rho_1 +\rho_2\right) 
\eeq

where $m_1$ and $m_2$ are, respectively, the masses of the light
and heavy fluids that move when the interface does. Substituting this
mass in \eq{eq:RT-Newton}

\beq
\ddot{\eta} = \frac{\left(\rho_2- \rho_1\right)}{\left(\rho_1 +\rho_2\right)} gk\eta 
\label{eq:RT-Newton2}
\eeq

which is an oscillation with frequency

\beq
\omega^2 = \frac{\left(\rho_2- \rho_1\right)}{\left(\rho_1 +\rho_2\right)} gk
\eeq

just as we've found via linear analysis with Fourier modes. The
physical arguments based on elementary hydrostatics used to
derive this equation show how the driving force $\rho_2- \rho_1 g \eta
A$ is created across the deformed interface.

\section{Kelvin-Helmholtz instability}

Consider a situation with no gravitational field, and with horizontal
shear between the fluids.

The equation for the boundary Langrangian perturbation is now 

\beq
\frac{\partial \eta}{\partial t} + \frac{\partial \eta}{\partial x} u_i = \frac{\partial \phi_i}{\partial z}. 
\label{eq:contactboundary-x}
\eeq

and the Bernoulli equation is

\beq
\rho_1\left(\frac{\partial \phi_1}{\partial t} +  \frac{1}{2}u^2\right) = \rho_2\left(\frac{\partial \phi_2}{\partial t} +  \frac{1}{2}u^2\right)
\label{eq:bernoulli-boundary-x1}
\eeq

expanding $u^2 = \bar{u}^2 + 2 \bar{u} u^\prime$ and writing
$u^\prime=\grad \phi$,

\beq
\rho_1\left(\frac{\partial \phi_1}{\partial t} +  u_1 \frac{\partial \phi_1}{\partial x}\right) = \rho_2\left(\frac{\partial \phi_2}{\partial t} +  u_2 \frac{\partial \phi_2}{\partial x}\right)
\label{eq:bernoulli-boundary-x2}
\eeq

expanding in Fourier modes, we find

\beqn
-i\left(\omega-k u_1\right) \hat{\eta} &=& -k \phi_1\\
-i\left(\omega-k u_2\right) \hat{\eta} &=&  k \phi_2\\
\rho_1 \hat{\phi}_1\left(\omega-k u_1\right) &=& \rho_2 \hat{\phi}_2\left(\omega-k u_2\right)
\eeqn

for which the dispersion relation is

\beq
\omega^2 \left(\rho_1+\rho_2\right) - 2\omega k
\left(\rho_1u_1+\rho_2u_2\right) + k^2 \left(\rho_1u_1^2+\rho_2u_2^2\right)
\eeq

The solution is

\beq
\omega_{1,2} = \frac{k\left(\rho_1u_1+\rho_2u_2\right)}{\rho_1+\rho_2}
\pm i \frac{\sqrt{\rho_1\rho_2}}{\rho_1+\rho_2}\vert u_1-u_2\vert
\eeq

Interestingly, in an ideal gas there is an imaginary growing mode component for every
$\vert u_1-u_2\vert > 0$, that is, {\it whenever there is shear}.

This means that a small wave-like perturbation at an interface will
grow rapidly into large waves that take the form of characteristic Kelvin-Helmholtz
``billows''. In the non-linear regime reached during the subsequent evolution of this
instability the waves are rolled up, leading to the creation of vortex like structures.
As the instability grows fastest for small scales (high k), with time the billows tend
to get larger and larger.

As the Kelvin-Helmholtz instability basically means that any sharp
velocity gradient in a shear flow is unstable in a freely streaming
fluid, this instability is particularly important for the creating of
fluid turbulence.

\section{Jeans Instability}

Lesson plan

\begin{itemize}
\item Intro (10 min)

  \begin{itemize}
    \item Jeans: star formation in the ISM, stucture formation in cosmology.
    \item Toomre: structure formation in galactic disks.
  \end{itemize}
  
\item Jeans (20 min)

  \begin{itemize}
  \item Physical intuition:Competition between pressure support and
    gravitational collapse.
    \item Hand-wavy: sound crossing time vs free fall. Sound removes
      compressibility. Gravity must act faster than that. Get Jeans
      length and mass.
    
    \item Mathematical setup: Linear perturbation on Continuity +
      Euler + Poisson. Compare Jeans less and mass obtained
      before. Jeans mass about 2 solar masses for molecular cloud conditions.

      \item Discussion: What kinds of perturbations grow or oscillate?
        Question: What happens if we increase the temperature? Or the density?
      
  \end{itemize}
  
\item Toomre (20 min)

  \begin{itemize}
    \item Physical picture:

        \begin{itemize}
        \item Why rotation changes everything.

        \item Differential rotation $\rightarrow$ shear $\rightarrow$ stabilization.

        \item Epicyclic motion \& angular momentum.
    \end{itemize}
    \item     Mathematical setup: Thin, axisymmetric disk with surface
      density $\Sigma$. Consider axisymmetric perturbations.

    \item Define epicyclic frequency.

    \item Toomre parameter

  \end{itemize}
  
\item Jeans vs Toomre (15 min)

  \begin{itemize}
    \item Jeans: infinite medium, 3D, pressure vs. gravity.
    \item Toomre: thin disk, 2D, includes shear and rotation.

     \item Plot \& interpret dispersion relation curves.

     \item Discussion: why disks are more stable than uniform media.

     \item Typical Q values in different astrophysical systems (e.g., Milky Way, protoplanetary disks).
  \end{itemize}
  
\item Wrap-up (10 min)

  \begin{itemize}
    \item Summary of key points: How dispersion relations encode the
      physical balance of forces. Critical scales and parameters (Jeans length, Toomre Q).

    \item Open questions in astrophysics: How do magnetic fields, turbulence, or cooling alter these instabilities?

  \end{itemize}
  
\end{itemize}



Jeans instability refers to the critical condition in which a cloud of
gas and dust, under the influence of its own gravity, begins to
collapse and form stars. This concept is pivotal in understanding star
formation, as it connects the initial mass function—the distribution
of masses for newly formed stars—to the rates at which these stars
form in molecular clouds. When the internal pressure of a cloud cannot
support the gravitational forces acting on it, the Jeans length
becomes crucial in determining whether collapse will occur.

The instability arises when the thermal pressure of a gas parcel within a cloud cannot
counterbalance its gravitational pull. This gas parcel eventually fragmenting and
forming stars. The fundamental question to ask is, what are the
wavelengths that go unstable, and what is the typical mass that collapses? 

The system is in hydrostatic equilibrium, and the perturbations evolve
according to the equations of hydrodynamics. The system is similar to
what we have considered so far, except that now, because the system is
selfgravitating, we need to consider Poisson's equation

\beqn
\frac{\partial \rho}{\partial t} + \left(\v{u}\cdot \del\right)\rho &=& -\rho\Div{\v{u}}\\
\frac{\partial \v{u}}{\partial t} + \left(\v{u}\cdot \del\right)\v{u}&=& -\frac{1}{\rho}\grad{p} - \grad{\varPhi}\\
\Laplace{\varPhi} &=& 4\pi G \rho
\eeqn

The system is closed with the equation of state $p=\rho c_s^2$. We
assume the system isothermal for simplicity.

Let us again decompose the system into base and perturbation

\beqn
\frac{\partial \rho^\prime}{\partial t} + \rho_0\Div{\v{u}^\prime} &=&0\\
\frac{\partial \v{u}^\prime}{\partial t} +\frac{c_s^2}{\rho_0}\grad{\rho^\prime} + \grad{\varPhi^\prime}&=&0\\
\Laplace{\varPhi^\prime} - 4\pi G \rho^\prime&=& 0
\eeqn

and apply a radial perturbation $\psi^\prime = \hat{\psi} \ {\rm exp} \left[i \left(kr - \omega t \right)\right]$

\beqn
-\omega \hat{\rho} +\rho_0k\hat{u}&=&0\\
-\omega\hat{u} +\frac{c_s^2}{\rho_0}k\hat{\rho} + k\hat{\varPhi}&=&0\\
k^2\hat{\varPhi} + 4\pi G \hat{\rho}&=&0
\eeqn

The system is 

\begin{equation}
\begin{vmatrix}
-\omega                       & \rho_0k    & 0 \\
\frac{c_s^2}{\rho_0}k  & -\omega  & k\\
 4\pi G & 0 & k^2
\end{vmatrix}	= 0 
\end{equation}

This leads to the dispersion relation

\beq
\omega^2  = k^2 c_s^2 - 4\pi G\rho_0 
\eeq

defining

\beq
k_J^2\equiv \frac{4\pi G\rho_0}{c_s^2} = \frac{4\pi G\rho_0 m_p \mu}{k_B T }
\eeq

we can write

\beq
\omega^2  = c_s^2 \left(k^2-k_J^2\right)
\eeq


Forr $k < k_J$ the complex eigenfrequency is imaginary and hence the
disturbance grows exponentially. The wavenumber $k_J$ defines the
{\it Jeans length}

\beq
\lambda_J \equiv \frac{2\pi}{k_J} = c_s\left(\frac{\pi}{G\rho_0}\right)^{1/2} = \left(\frac{\pi k_B T }{G\rho_0 m_p \mu}\right)^{1/2}
\eeq

It also defines the {\it Jeans mass}

\beq
M_J \equiv  \left(\frac{\lambda_J}{2}\right)^3 \rho_0 = \left(\frac{\pi k_B T }{4G m_p  \mu}\right)^{3/2} \frac{1}{\rho_0^{1/2}}
\eeq

Perturbations of size larger than this in a gas cloud would grow,
become self-gravitating and collapse. For $T=10$\,K, $\mu= 2.3$, and
number density $n = 10^3$ cm$^{-3}$, the Jeans mass is about 2 solar
masses. 

Notice that the dispersion relation $\omega^2  =  k^2 c_s^2 - 4\pi
G\rho_0  $ evidences the physics of the problem. Gravity and pressure
are competing. Gravity is trying to compress the gas, sounds waves
remove the compressibility. Sound waves remove the compressibility at
the sound crossing time $R/c_s$. Gravity compresses at the free fall
time $1/\sqrt{G\rho}$. There exists a length $R$ for which they're
equal. For larger lengths, the sound crossing time is longer than the
free fall time, and the compressibility cannot be removed. The system
collapses. This length is the Jeans length

\beq
\lambda_J \sim \frac{c_s}{\sqrt{G\rho}}
\eeq

which is what we derived from rigorous linear analysis except for a
factor $\sqrt{\pi}$.


In the early universe, this applies to baryonic matter after recombination (when it decoupled from radiation). Before that, photon pressure prevented collapse. After recombination ($\sim$380,000 years after the Big Bang), perturbations in dark matter and baryons could grow under Jeans instability, seeding galaxies.


Some numbers: for molecular cloud

Let us evaluate the Jeans length and mass, Plugging in numbers typical of dense molecular cores 

\beqn
\lambda_J &\approx& 1.0\,pc \left( \frac{T}{10 \ {\rm K}}\right)^{1/2} \left( \frac{n}{10^3 \ {\rm cm}^{-3}}\right)^{1/2}\\
M_J &\approx& 26\,M_\sun \left( \frac{T}{10 \ {\rm K}}\right)^{3/2} \left( \frac{n}{10^3 \ {\rm cm}^{-3}}\right)^{1/2}
\eeqn

\beq
t_ff = {G\rho}^{-1/2} \approx 2.1 Myr \left( \frac{n}{10^3 \ {\rm cm^{-3}}}\right)^{-1/2}
\eeq

\section{Toomre Instability}

Jeans instability is responsible for structure formation in
homogeneous media. Yet, we also observe structure formation in 
flattened disks, be that Galaxies or circumstellar and circumplanetary
disks. In this case, fragmentation is due to {\it Toomre} instability. 

In these systems, the mostly 2D configuration and
the presence of differential rotation (shear) will
lead to different stability properties. In addition to sound waves removing compressibility, in the presence of differential rotation, a collapsing region
will also now have to contend with shear. The collapse has to be
faster than the timescale in which the shear is advecting the blobs
apart.

To collapse, a region must contract gravitationally faster than acoustics can remove the compression and differential
rotation can shear it away. We can construct a criteron based on these
constraints, as we did for Jeans instability 

The timescale for acoustics is the sound-crossing time

\beq
t_c = \lambda/c_s
\eeq

The timescale for gravity is the free-fall time

\beq
t_g = \frac{1}{\sqrt{G\rho}}
\eeq

The timescale for shear is the orbital time

\beq
t_s = \varOmega^{-1}
\eeq

Constructing a dimensionless number from ``Gravity vs Shear and
Pressure'', we find

\beqn
Q &\sim& \frac{t_g^2}{t_c t_s}\\
&\sim&\frac{c_s\varOmega}{G\rho}
\eeqn

If $Q >$ 1, gravity is slower than the combined action of pressure and shear, and the system is stable. If $Q<1$, gravity is faster, and the
system collapses. This analysis is off from the actual $Q$ from rigorous linear analysis by a factor $\pi^{-1}$, and the fact that shear stability is
given by the epicyclic frequency $\kappa$ instead. The dispersion
relation is

\beq
\omega^2 = \kappa^2 = 2\pi G\Sigma|k| + c_s^2k^2
\eeq

Gravity pushes $\omega$ to imaginary values. Pressure and rotation push it to real values.

In the saturated state, it is seen that an adiabatic system evolves
toward $Q\sim 1$. If the system compresses (lowering $Q$), it also
heats up (increasing $Q$). If the system expands (increasing $Q$), it
also cools down (decreasing $Q$). The process is self-regulated,
leading to gravitational turbulence and spiral arms.
Cooling has to be fast to allow for collapse.

Toomre instability is slow cooling disks as the disk mass is increased.
The system passes from gravitational turbulence (low mass) to spiral arms (higher mass).
If the mass is further increased, collapse is inevitable. Toomre instability might be responsible for the
formation of giant planets in wide orbits. The process is also called simply “gravitational
instability” or “disk instability” in the planet formation literature.


